\documentclass[troisieme]{fiche}
\metadonnees{8}{Humbug}{Bloc B — Ce qu'on peut, ce qu'on doit}

\begin{document}
\entete

\objectifs{%
\textbf{Langue} : \emph{must}, \emph{mustn't}, \emph{have to} ; traduire \og il faut \fg.\par
\textbf{Texte} : Charles Dickens, \emph{A Christmas Carol} (1843), premier chant.\par
\textbf{Durée} : 1 h — exercices 20 min, texte et questions 25 min, version 15 min.}

%% ===================================================================
\exercice[12 min]{\emph{Must}, \emph{mustn't}, \emph{have to}}

\rappel{\textbf{Rappel.} \emph{Must} vient de celui qui parle, \emph{have to} d'une règle
extérieure. \emph{Mustn't} interdit ; \emph{don't have to} dispense.}

\consigne{Complétez.}

\begin{anglais}
\begin{enumerate}[label=\arabic*.,itemsep=2.5mm,leftmargin=*]
  \item The clerk \trou{has to} work in a cold little room: Scrooge keeps the coal.
  \item You \trou{mustn't} touch the coal-box without asking.
  \item ``I \trou{must} finish these letters before six,'' said the clerk.
  \item You \trou{don't have to} come on Christmas Day; the office is closed.
  \item In England, children \trou{have to} go to school until sixteen.
  \item We \trou{mustn't} be late for dinner.
  \item Scrooge \trou{had to} pay his clerk, but he paid him very little.
  \item You \trou{don't have to} shout: I can hear you perfectly.
  \item ``You \trou{must} come and dine with us tomorrow,'' said the nephew.
  \item Visitors \trou{mustn't} enter the counting-house without knocking.
\end{enumerate}
\end{anglais}

\note{\emph{Must} : 3, 9. \emph{Have to} : 1, 5, 7. \emph{Mustn't} : 2, 6, 10.
\emph{Don't have to} : 4, 8.}

\note[Le piège]{\emph{Mustn't} n'est pas la négation de \emph{must}. \emph{You mustn't
go} = il ne faut pas y aller. \emph{You don't have to go} = ce n'est pas la peine d'y aller.
Et \emph{must} n'a pas de passé : on emprunte \emph{had to} (7).}

%% ===================================================================
\exercice[8 min]{Traduire \og il faut \fg}

\rappel{\textbf{Rappel.} Le français dit \og il faut \fg{} pour tout ; l'anglais choisit
entre \emph{must}, \emph{have to}, \emph{need} et \emph{it takes}.}

\consigne{Traduisez en anglais.}

\begin{anglais}
\begin{enumerate}[label=\arabic*.,itemsep=3mm,leftmargin=*]
  \item Il faut que tu partes maintenant.
    \lignes{1}
    \corr{You must go now. \emph{ou} You have to go now.}
  \item Il ne faut pas rire.
    \lignes{1}
    \corr{You mustn't laugh.}
  \item Il faut trois heures pour aller à Londres.
    \lignes{1}
    \corr{It takes three hours to get to London.}
  \item Il me faut un nouveau manteau.
    \lignes{1}
    \corr{I need a new coat.}
  \item Il a fallu qu'il travaille le jour de Noël.
    \lignes{1}
    \corr{He had to work on Christmas Day.}
  \item Il ne faut pas que vous soyez en retard.
    \lignes{1}
    \corr{You mustn't be late.}
\end{enumerate}
\end{anglais}

\note[Quatre traductions pour un seul verbe]{Obligation : \emph{must} ou \emph{have to}.
Interdiction : \emph{mustn't}. Besoin d'une chose : \emph{need}. Durée nécessaire :
\emph{it takes}. Le français ne distingue pas, l'anglais oblige à choisir — d'où l'intérêt de
l'exercice.}

%% ===================================================================
\newpage
\section{Texte}

\noindent\small\textit{La veille de Noël, dans le bureau de Scrooge, un vieil usurier. Son
employé copie des lettres dans un réduit glacé.}\normalsize

\begin{texte}
The door of Scrooge's \gl[a counting-house]{counting-house}{un cabinet comptable} was open
that he might keep his eye upon his \gl[a clerk]{clerk}{un employé aux écritures}, who in a
\gl[dismal]{dismal}{lugubre} little cell beyond, a sort of tank, was copying letters. Scrooge
had a very small fire, but the clerk's fire was so very much smaller that it looked like one
\gl[coal]{coal}{du charbon}. But he couldn't \gl[to replenish]{replenish}{regarnir} it, for
Scrooge kept the coal-box in his own room; and so surely as the clerk came in with the
\gl[a shovel]{shovel}{une pelle}, the master predicted that it would be necessary for them to
part.

``A merry Christmas, uncle! God save you!'' cried a cheerful voice. It was the voice of
Scrooge's nephew, who came upon him so quickly that this was the first
\gl[an intimation]{intimation}{un signe, un avertissement} he had of his approach.

``Bah!'' said Scrooge, ``\gl[humbug]{Humbug}{fariboles, sottises}!''

He had so heated himself with rapid walking in the fog and frost, this nephew of Scrooge's,
that he was all in a \gl[a glow]{glow}{un éclat, une chaleur}; his face was
\gl[ruddy]{ruddy}{coloré} and handsome; his eyes sparkled, and his breath smoked again.

``Christmas a humbug, uncle!'' said Scrooge's nephew. ``You don't mean that, I am sure?''

``I do,'' said Scrooge. ``Merry Christmas! What right have you to be merry? What reason have
you to be merry? You're poor enough.''

``Come, then,'' returned the nephew \gl[gaily]{gaily}{gaiement}. ``What right have you to be
\gl[dismal]{dismal}{sombre}? What reason have you to be \gl[morose]{morose}{morose}? You're
rich enough.''

Scrooge having no better answer ready on the \gl[on the spur of the moment]{spur of the
moment}{sur le coup}, said, ``Bah!'' again; and followed it up with ``Humbug.''

``Don't be cross, uncle!'' said the nephew.

``What else can I be,'' returned the uncle, ``when I live in such a world of fools as this?
Merry Christmas! What's Christmas time to you but a time for paying \gl[a bill]{bills}{une
facture} without money; a time for finding yourself a year older, but not an hour richer? If
I could work my will,'' said Scrooge indignantly, ``every idiot who goes about with `Merry
Christmas' on his lips should be boiled with his own pudding, and buried with a
\gl[a stake]{stake}{un pieu} of \gl[holly]{holly}{du houx} through his heart. He should!''

``Uncle!'' \gl[to plead]{pleaded}{supplier} the nephew.

``Nephew!'' returned the uncle \gl[sternly]{sternly}{sévèrement}, ``keep Christmas in your own
way, and let me keep it in mine.''

``Keep it!'' repeated Scrooge's nephew. ``But you don't keep it.''

``Let me leave it alone, then,'' said Scrooge.
\end{texte}
\source{Charles Dickens, \emph{A Christmas Carol}, Londres, Chapman \& Hall, 1843, premier
chant.}

\partiel[15 min]{Questions}
\consigne{Répondez aux deux premières questions en français, aux deux dernières en anglais.}

\begin{enumerate}[label=\arabic*.,itemsep=2.5mm,leftmargin=*]
  \item Comment Scrooge traite-t-il son employé ? Relevez trois détails.
    \lignes{3}
    \corr{Il laisse sa porte ouverte pour le surveiller ; l'employé travaille dans un réduit
    lugubre ; son feu est si petit qu'il ressemble à un seul morceau de charbon ; Scrooge
    garde la boîte à charbon dans son bureau et menace de le renvoyer dès qu'il vient avec la
    pelle.}
  \item Sur quoi porte exactement la dispute entre l'oncle et le neveu ?
    \lignes{2}
    \corr{Sur Noël. Le neveu le souhaite joyeux, Scrooge répond que c'est une sottise : pour
    lui, c'est le moment de payer des factures sans argent et de se découvrir plus vieux d'un
    an sans être plus riche d'une heure.}
  \item \en{The nephew answers Scrooge with Scrooge's own words. Find that answer and
    explain it.}
    \lignes{3}
    \corr{Scrooge asks \emph{What right have you to be merry? You're poor enough.} The nephew
    gives the sentence back: \emph{What right have you to be dismal? You're rich enough.} He
    shows that Scrooge's rule works both ways, and Scrooge has no answer but \emph{Bah!}}
  \item \en{Find in the text one sentence with \emph{couldn't} and one with \emph{should}.
    What does each one tell us?}
    \lignes{3}
    \corr{\emph{But he couldn't replenish it} : le clerc n'en a pas la possibilité, parce que
    Scrooge garde le charbon. \emph{Every idiot\ldots should be boiled with his own pudding} :
    ce n'est pas un conseil mais un souhait féroce — ce que Scrooge ferait s'il pouvait.}
\end{enumerate}

%% ===================================================================
\section{Version}
\consigne{Traduisez le premier paragraphe du texte, de \emph{The door of Scrooge's
counting-house} à \emph{necessary for them to part}.}

\lignes{12}

\corr{\textbf{Proposition de traduction.} La porte du cabinet de Scrooge restait ouverte, afin
qu'il pût garder l'œil sur son employé, lequel, dans un petit réduit lugubre, une sorte de
citerne, copiait des lettres. Scrooge avait un tout petit feu, mais celui de l'employé était
si nettement plus petit encore qu'il avait l'air d'un unique morceau de charbon. Impossible
pour lui de le regarnir, car Scrooge gardait la boîte à charbon dans sa propre pièce ; et
aussi sûrement que l'employé entrait avec la pelle, le maître annonçait qu'il leur faudrait se
séparer.}

\note[Choix de traduction 1]{\emph{that he might keep his eye upon his clerk} : \emph{that} +
\emph{might} exprime le but, \og afin qu'il pût \fg. Ce n'est pas une possibilité mais une
intention.}

\note[Choix de traduction 2]{\emph{it would be necessary for them to part} : le style de
Scrooge est celui d'un homme qui n'annonce jamais un renvoi directement. Garder la lourdeur
(\og qu'il leur faudrait se séparer \fg) : elle fait le personnage.}

\note[Choix de traduction 3]{\emph{so surely as} = \og aussi sûrement que \fg. La phrase dit
que la menace revenait à chaque fois, mécaniquement : c'est un imparfait d'habitude en
français.}

%% ===================================================================
\partiel{Lexique à mémoriser}
\begin{multicols}{2}\small
\begin{itemize}[label=--,itemsep=0.8mm,leftmargin=*]
  \item \textbf{a clerk} — un employé de bureau
  \item \textbf{coal} — le charbon
  \item \textbf{a fire} — un feu
  \item \textbf{dismal} — lugubre, sombre
  \item \textbf{cheerful} — enjoué, gai
  \item \textbf{a nephew} — un neveu ; \textit{féminin} a niece
  \item \textbf{merry} — joyeux \textit{(fiche 2)}
  \item \textbf{fog} — le brouillard ; \textit{adjectif} foggy
  \item \textbf{frost} — le gel
  \item \textbf{cross} — de mauvaise humeur
  \item \textbf{a fool} — un idiot
  \item \textbf{a bill} — une facture, une note
\end{itemize}
\end{multicols}

\end{document}
